\subsection*{Discussion: conditional entropy and group actions}

The entropy viewpoint is useful even before tensor terminology enters:
the coset label specifies conditioning data, and
$\HH(X\mid X\bmod V)$ measures the entropy left within those fibers.
The question in Lemma~\ref{lem:invariant} is how much of the total
entropy can lie in a subspace of a given dimension when the law has
irreducible linear symmetry. Its application to Dicke tensors turns
this finite entropy comparison into the precise exponential rate of
extractable diagonal structure for every multiplicity tuple.

Related dynamical subjects include relative sequence entropy for
amenable group actions \cite{LWW2025} and sequence entropy and IT-tuples
for minimal group actions \cite{LWX2025}. These are background topics
for comparison; no theorem from either work is invoked in this paper.
We make no identification of sequence entropy with the finite Shannon
entropy used here, or of IT-tuples with the induced matchings used to
extract a diagonal tensor. The hypotheses of our entropy lemma remain
finite support, a positive finite-dimensional rational vector space,
and an invariant law under an irreducible rational-linear action.

\paragraph{Open questions.}
The all-multiplicity asymptotic-subrank formula is settled by
Theorem~\ref{thm:main}. Exact finite-power subranks, convergence rates
uniform over multiplicities and efficient extraction of the restrictions
are not determined here. A further question for comparison is whether
quotient-entropy estimates can guide combinatorial selection in related
dynamical settings. No such extension is asserted.
